\scanpage{112}
\begin{center}- \mspage{123} -\end{center}
\margin{IX}
et par passage aux modèles de Néron définit donc un homomorphisme sur les tores maximaux des fibres spéciales, donc en sens inverse un homomorphisme sur les Groupes de caractères, d'où un homomorphisme
\[
(10.2.6)\qquad \xi_t:\overline M\longrightarrow\underline M'\quad,
\]
où $\overline M$ est défini en termes de $\overline A_K$ comme $\underline M$ en termes de $A_K$. On peut donc former l'accouplement
\[
(10.2.7)\qquad \varphi_\xi^t:\underline M_\ell\otimes\overline{\underline M}_\ell\longrightarrow\mathbb Z_\ell\quad,\quad
\varphi^t(x,y)=\varphi^t(x,\xi_t(y))\quad.
\]
Notons en passant que pour un schéma abélien $A_K$ à réduction semi-stable variable, le Groupe constant tordu $\underline M$ est un foncteur covariant en $A_K$, et qu'alors la formation des accouplements (10.2.7) associés aux biextensions de schémas abéliens à réduction semi-stable sur $(K,S)$ est compatible, en un sens évident, avec la formation des images inverses de biextensions. Bien entendu, lorsque $\xi$ est la biextension canonique (``de Weil'') de $(A_K,A'_K)$, l'accouplement (10.2.7) n'est autre que l'accouplement de monodromie (9.2.1). - Ceci posé, la relation de symétrie peut s'exprimer de la façon suivante : considérons la symétrique ${}^s\xi$ de $\xi$, qui est une biextension de $(\overline A_K,A_K)$, d'où un accouplement
\[
\varphi_{{}^s\xi}^t:\overline{\underline M}_\ell\otimes\underline M_\ell\longrightarrow\mathbb Z_\ell\quad;
\]
on a alors la relation de symétrie
\[
(10.2.8)\qquad\varphi_{{}^s\xi}^t=\varphi_\xi^t\circ s\quad,\quad
\text{où}\quad s:\overline{\underline M}_\ell\otimes\underline M_\ell\xrightarrow{\sim}\underline M_\ell\otimes\overline{\underline M}_\ell
\]
est encore l'isomorphisme de symétrie. On comparera cette relation à la relation d'antisymétrie de VIII 2.2.11.
\margin{441}
\margin{441}
\margin{111}
